\chapter{Linear Algebra}\label{ch:b2:linalg}

The linear algebra of the Year 1 volume worked over $\R$ or $\C$ in
finite dimension, and admitted the general determinant. This chapter
upgrades all three restrictions: the theory is stated over an
arbitrary field $K$, the interplay between a space and its
\emph{dual} is developed systematically (dual bases, annihilators,
transposes), and the determinant is finally \emph{constructed} from
alternating multilinear forms and the signature of
\cref{ch:b2:structures} --- discharging every admission of Year 1.

Throughout, $K$ is a field ($\Q$, $\R$, $\C$, or $\Z/p\Z$ --- the
theory does not care) and, unless stated, spaces are
finite-dimensional over $K$. The Year 1 results (bases, dimension,
rank--nullity, matrices) transfer verbatim: their proofs never used
anything but the field axioms.

\section{Dual space}

\begin{definition}[Dual space, dual basis]\label{def:b2:linalg:dual}
The \emph{dual}\index{dual space} of $E$ is $E^* = \mathcal{L}(E,
K)$, the space of linear forms. If $\mathcal{B} = (e_1, \dots, e_n)$
is a basis of $E$, the \emph{coordinate forms} $e_1^*, \dots, e_n^*$
defined by $e_i^*(e_j) = \delta_{ij}$ (Kronecker: $1$ if $i = j$,
else $0$) form the \emph{dual basis}\index{dual basis}
$\mathcal{B}^*$ of $E^*$; in particular $\dim E^* = \dim E$, and
\[
x = \sum_{i=1}^{n} e_i^*(x)\, e_i \quad (x \in E),
\qquad
\varphi = \sum_{i=1}^{n} \varphi(e_i)\, e_i^* \quad (\varphi \in
E^*).
\]
\end{definition}

\begin{proof}[Proof that $\mathcal{B}^*$ is a basis]
Free: applying a null combination $\sum \lambda_i e_i^* = 0$ to
$e_j$ gives $\lambda_j = 0$. Generating: for $\varphi \in E^*$, the
form $\varphi - \sum_i \varphi(e_i) e_i^*$ kills every $e_j$, hence
is zero (a linear map vanishing on a basis vanishes). The two
display formulas are the same computations read forwards.
\end{proof}

\begin{example}\label{ex:b2:linalg:dualexamples}
On $K_n[X]$ with basis $(1, X, \dots, X^n)$: the dual basis is $P
\mapsto \frac{P^{(k)}(0)}{k!}$ (Taylor coefficients). Another basis
of the dual: the evaluations $P \mapsto P(x_i)$ at $n + 1$ distinct
points --- its ``pre-dual'' basis in $K_n[X]$ is exactly the family
of Lagrange polynomials $L_i$ (Year 1 volume), since $L_i(x_j) =
\delta_{ij}$. Interpolation \emph{is} duality.
\end{example}

\begin{method}[Dual and antedual bases in practice]\label{met:b2:linalg:dualbasis}
To expand a form $\varphi$ on a basis $(e_i)$ of $E$: the
coordinates are the \emph{values} $\varphi(e_i)$ --- no system to
solve. To find the basis $(u_j)$ of $E$ whose dual is a given
basis $(\varphi_1, \dots, \varphi_n)$ of $E^*$ (the
\emph{antedual}): solve the $n$ linear systems
\[
\varphi_i(u_j) = \delta_{ij} \qquad (1 \leq i \leq n),
\]
one column $u_j$ at a time; in matrix terms, if the rows of $M$
list the coefficients of the $\varphi_i$ in a known basis of
$E^*$, the columns of $M^{-1}$ are the $u_j$. Existence and
uniqueness of the antedual are proved in this chapter's weekend
problem; the computation is always this inversion.
\end{method}

\begin{example}[A dual basis of $\R^2$, fully computed]\label{ex:b2:linalg:dualbasis2d}
For the basis $b_1 = (1, 1)$, $b_2 = (1, -1)$ of $\R^2$: the
dual basis $(b_1^*, b_2^*)$ must satisfy $b_i^*(b_j) =
\delta_{ij}$. Writing $b_1^*(x, y) = \alpha x + \beta y$, the
conditions $\alpha + \beta = 1$ and $\alpha - \beta = 0$ give
\[
b_1^*(x, y) = \frac{x + y}{2},
\qquad\text{and likewise}\qquad
b_2^*(x, y) = \frac{x - y}{2} .
\]
Sanity checks: $b_1^*$ is \emph{not} $e_1^* + e_2^*$ evaluated
naively --- the dual basis depends on the whole basis, not on
each vector separately (replacing $b_2$ by $(0, 1)$ changes
$b_1^*$ into $x \mapsto x$). And the expansion formula works:
$(x, y) = \frac{x+y}2\,b_1 + \frac{x-y}2\,b_2$, the even/odd
decomposition of a pair --- dual bases are coordinate
extractors, and this one extracts symmetric and antisymmetric
parts.
\end{example}

\begin{definition}[Annihilator]\label{def:b2:linalg:annihilator}
For a subspace $F \subseteq E$, the
\emph{annihilator}\index{annihilator} is
\[
F^{\circ} = \{\varphi \in E^* : \varphi|_F = 0\},
\]
a subspace of $E^*$.
\end{definition}

\begin{theorem}[Dimension of the annihilator]\label{thm:b2:linalg:annihilator}
$\dim F^{\circ} = \dim E - \dim F$. Moreover $F \mapsto F^\circ$
reverses inclusions, and $F$ is recovered from its annihilator:
\[
F = \{x \in E : \forall\varphi \in F^\circ,\ \varphi(x) = 0\}.
\]
Consequently every subspace of dimension $p$ in dimension $n$ is the
solution set of $n - p$ independent linear equations --- and
conversely.
\end{theorem}

\begin{proof}
Choose a basis $(e_1, \dots, e_p)$ of $F$ completed into a basis of
$E$. A form $\varphi = \sum \varphi(e_i) e_i^*$ annihilates $F$ iff
its first $p$ coefficients vanish: $F^\circ =
\operatorname{Vect}(e_{p+1}^*, \dots, e_n^*)$, of dimension $n - p$.
Inclusion reversal is immediate. For the recovery: the right-hand
side contains $F$; conversely, if $x \notin F$, complete a basis of
$F$ by $x$ and further vectors; the coordinate form of $x$ in this
basis annihilates $F$ but not $x$. The ``equations'' reading takes a
basis $(\varphi_1, \dots, \varphi_{n-p})$ of $F^\circ$: then $F =
\bigcap \ker\varphi_j$, an intersection of $n - p$ independent
hyperplanes.
\end{proof}

\begin{example}[An annihilator, both directions]\label{ex:b2:linalg:annihilatorwork}
Let $F = \operatorname{Vect}\bigl((1, 2, 1),\ (1, 0, -1)\bigr)
\subseteq \R^3$. A form $\varphi = a\,e_1^* + b\,e_2^* + c\,e_3^*$
annihilates $F$ iff
\[
a + 2b + c = 0
\qquad\text{and}\qquad
a - c = 0 ,
\]
i.e.\ $c = a$ and $b = -a$: $F^\circ = \R\,(e_1^* - e_2^* +
e_3^*)$, of dimension $3 - 2 = 1$ as
\cref{thm:b2:linalg:annihilator} requires. Reading it backwards:
$F = \{(x, y, z) : x - y + z = 0\}$ --- the plane recovered as
the kernel of the single form spanning $F^\circ$. Going from a
spanning family to equations \emph{is} computing an annihilator;
going from equations to a parametrization is computing a
pre-annihilator. (Check: both spanning vectors satisfy $x - y +
z = 0$.)
\end{example}

\begin{definition}[Transpose map]\label{def:b2:linalg:transpose}
For $u \in \mathcal{L}(E, F)$, the \emph{transpose}\index{transpose
map} $u^{\mathsf T} \in \mathcal{L}(F^*, E^*)$ is
\[
u^{\mathsf T}(\psi) = \psi \circ u .
\]
It satisfies $(v \circ u)^{\mathsf T} = u^{\mathsf T} \circ v^{\mathsf
T}$, and in dual bases, the matrix of $u^{\mathsf T}$ is the
transposed matrix of $u$ --- which finally \emph{explains} the
transpose of Year 1.
\end{definition}

\begin{example}[The transpose, entry by entry]\label{ex:b2:linalg:transposework}
Let $u \colon \R^2 \to \R^3$ have matrix $A =
\begin{psmallmatrix} 1 & 2\\ 0 & 1\\ 3 & 0\end{psmallmatrix}$ in
the canonical bases. For $\psi = b_1f_1^* + b_2f_2^* + b_3f_3^*
\in (\R^3)^*$, compute $u^{\mathsf T}(\psi) = \psi \circ u$ on
the basis of $\R^2$:
\[
(\psi \circ u)(e_1) = \psi(1, 0, 3) = b_1 + 3b_3,
\qquad
(\psi \circ u)(e_2) = \psi(2, 1, 0) = 2b_1 + b_2 .
\]
So $u^{\mathsf T}(\psi) = (b_1 + 3b_3)\,e_1^* + (2b_1 +
b_2)\,e_2^*$, and in the dual bases the matrix of $u^{\mathsf T}$
is
\[
\begin{pmatrix} 1 & 0 & 3\\ 2 & 1 & 0\end{pmatrix}
= A^{\mathsf T} :
\]
the abstract transpose \emph{is} the flipped matrix, with no
computation left to believe on faith. Note the mechanism: the
$j$-th \emph{column} of $A$ became the $j$-th \emph{row} of the
new matrix because $\psi \circ u$ reads $u$'s outputs through
$\psi$'s coefficients.
\end{example}

\begin{proposition}\label{prop:b2:linalg:transposerank}
$\ker u^{\mathsf T} = (\operatorname{im} u)^{\circ}$ and
$\operatorname{im} u^{\mathsf T} = (\ker u)^{\circ}$. Consequently
$\operatorname{rk}(u^{\mathsf T}) = \operatorname{rk}(u)$: row rank
equals column rank, proved structurally.
\end{proposition}

\begin{proof}
$\psi \in \ker u^{\mathsf T} \iff \psi \circ u = 0 \iff \psi$ kills
$\operatorname{im} u$: the first identity. For the second: $u^{\mathsf
T}(\psi) = \psi \circ u$ kills $\ker u$ always, so
$\operatorname{im} u^{\mathsf T} \subseteq (\ker u)^\circ$;
dimensions match by rank--nullity and
\cref{thm:b2:linalg:annihilator}:
\[
\operatorname{rk} u^{\mathsf T} = \dim F^* - \dim\ker u^{\mathsf T}
= \dim F - \bigl(\dim F - \operatorname{rk} u\bigr)
= \operatorname{rk} u
= \dim (\ker u)^{\circ} . \qedhere
\]
\end{proof}

\begin{example}[Rank read on both sides]\label{ex:b2:linalg:rankbothsides}
Let
\[
A = \begin{pmatrix}
1 & 2 & 0 & 1\\
0 & 1 & 1 & 1\\
1 & 3 & 1 & 2
\end{pmatrix} .
\]
\emph{Column rank:} the third row is the sum of the first two, so
$\operatorname{rk} A \leq 2$; columns $1$ and $2$ are free:
$\operatorname{rk} A = 2$. \emph{The transpose's kernel:} solving
$A^{\mathsf T}y = 0$ gives $y \in \R\,(1, 1, -1)$, so $\ker
A^{\mathsf T}$ has dimension $1 = 3 - 2$: exactly
$(\operatorname{im} A)^\circ$ under the identification of
$(\R^3)^*$ with row vectors, as
\cref{prop:b2:linalg:transposerank} asserts --- the single
relation ``row$_3$ = row$_1$ + row$_2$'' \emph{is} the
annihilator of the column space. Row rank ($2$ free rows) and
column rank agree not by accident but because both equal
$\operatorname{rk} A = \operatorname{rk} A^{\mathsf T}$.
\end{example}

\begin{example}[Duality reads a quadrature rule]\label{ex:b2:linalg:quadratureread}
Why does a rule like Simpson's (\cref{exo:b2:linalg:4}) exist
and why is it unique? Duality answers before any computation.
On $E = \R_2[X]$, the integral $P \mapsto \int_0^1 P$ is one
specific vector of the three-dimensional dual $E^*$; the
evaluations at $0$, $\frac12$, $1$ form a \emph{basis} of
$E^*$; hence the integral expands uniquely on them --- that
expansion \emph{is} Simpson's rule, coefficients included. A
dimension count also calibrates expectations: on $\R_3[X]$,
four dimensions of forms cannot in general be spanned by three
evaluations, so exactness on cubics is not owed by duality;
that Simpson integrates cubics exactly anyway is a bonus
symmetry (odd-degree cancellation around $\frac12$), to be
checked by hand. Rules with $n + 1$ nodes are expansions of the
integration form in an evaluation basis of $\R_n[X]^*$:
existence and uniqueness cost one dual-basis theorem; only the
bonus degrees cost work.
\end{example}

\section{Multilinear alternating forms}

\begin{definition}\label{def:b2:linalg:alternating}
A map $f \colon E^n \to K$ is \emph{$n$-linear} when it is linear in
each variable, and \emph{alternating}\index{alternating form} when
it vanishes whenever two arguments are equal. Alternating implies
\emph{antisymmetric}: swapping two arguments changes the sign
(expand $f(\dots, x + y, \dots, x + y, \dots) = 0$); more generally,
for $\sigma \in \mathfrak{S}_n$,
\[
f(x_{\sigma(1)}, \dots, x_{\sigma(n)}) =
\varepsilon(\sigma)\, f(x_1, \dots, x_n),
\]
by decomposing $\sigma$ into transpositions
(\cref{thm:b2:structures:signature}).
\end{definition}

\begin{theorem}[The fundamental theorem of determinants]\label{thm:b2:linalg:detspace}
Let $\dim E = n$ and $\mathcal{B} = (e_1, \dots, e_n)$ a basis. The
space of alternating $n$-linear forms on $E$ has dimension $1$:
every such form is a multiple of
\[
\det{}_{\mathcal{B}}(x_1, \dots, x_n)
= \sum_{\sigma \in \mathfrak{S}_n} \varepsilon(\sigma)
\prod_{i=1}^{n} a_{\sigma(i),\,i},
\qquad
x_j = \sum_{i} a_{ij} e_i ,
\]
and $\det_{\mathcal{B}}$ is the unique one taking the value $1$ on
$\mathcal{B}$.\index{determinant!construction}
\end{theorem}

\begin{proof}
Let $f$ be alternating $n$-linear. Expanding each argument on
$\mathcal{B}$ by multilinearity,
\[
f(x_1, \dots, x_n)
= \sum_{i_1, \dots, i_n} a_{i_1,1}\cdots a_{i_n,n}\,
f(e_{i_1}, \dots, e_{i_n}).
\]
Terms with a repeated index vanish (alternating); the surviving
tuples $(i_1, \dots, i_n)$ are the injective ones, i.e.\ $i_k =
\sigma(k)$ for a permutation $\sigma$, and antisymmetry reorders
$f(e_{\sigma(1)}, \dots, e_{\sigma(n)}) = \varepsilon(\sigma)
f(e_1, \dots, e_n)$. Hence
\[
f = f(e_1, \dots, e_n) \cdot \det{}_{\mathcal{B}} :
\]
every alternating form is that multiple, provided
$\det_{\mathcal{B}}$ itself (the displayed sum) \emph{is}
alternating $n$-linear and takes value $1$ on $\mathcal B$.
Multilinearity is clear (each summand is linear in each column).
Value on $\mathcal B$: the only nonzero term is $\sigma =
\mathrm{id}$. Alternating: suppose $x_j = x_k$ ($j \neq k$), so
that the coordinate columns satisfy $a_{i j} = a_{i k}$ for all
$i$. Pair each $\sigma$ with $\sigma' = \sigma\circ(j\,k)$ ---
an involution without fixed points on $\mathfrak{S}_n$. The
paired products coincide:
\[
\prod_i a_{\sigma'(i),\,i}
= a_{\sigma(k),\,j}\; a_{\sigma(j),\,k}
\prod_{i \neq j,k} a_{\sigma(i),\,i}
= a_{\sigma(k),\,k}\; a_{\sigma(j),\,j}
\prod_{i \neq j,k} a_{\sigma(i),\,i}
= \prod_i a_{\sigma(i),\,i},
\]
using the equality of the columns $j$ and $k$; while
$\varepsilon(\sigma') = -\varepsilon(\sigma)$. Each pair
contributes zero: the sum vanishes.
\end{proof}

\begin{example}[Sarrus, derived and demolished]\label{ex:b2:linalg:sarrus}
For $n = 3$ the permutation formula has exactly $3! = 6$ terms.
Listing $\mathfrak{S}_3$ by signature ---
$\mathrm{id}$, $(1\,2\,3)$, $(1\,3\,2)$ even;
$(1\,2)$, $(1\,3)$, $(2\,3)$ odd --- gives
\[
\det A = a_{11}a_{22}a_{33} + a_{21}a_{32}a_{13} +
a_{31}a_{12}a_{23}
- a_{21}a_{12}a_{33} - a_{31}a_{22}a_{13} -
a_{11}a_{32}a_{23} :
\]
precisely the ``diagonals'' rule of Sarrus taught in school ---
now a theorem, with the mysterious signs identified as
signatures. The demolition: for $n = 4$ there are $24$
permutations, of which only $8$ are picked up by any
diagonal-drawing scheme; Sarrus has no degree-$4$ version, and
cofactor expansion (\cref{thm:b2:linalg:detrules}~(4)) takes
over. Counting terms is also a warning: the permutation formula
has $n!$ summands, so it is a \emph{definition}, not an
algorithm --- row reduction computes $\det$ in $O(n^3)$
operations instead.
\end{example}

\begin{definition}[Determinants]\label{def:b2:linalg:det}
The \emph{determinant of a family} in a basis is
$\det_{\mathcal{B}}(x_1, \dots, x_n)$; the \emph{determinant of a
matrix} $A$ is the determinant of its columns in the canonical basis
--- the permutation formula above; the \emph{determinant of an
endomorphism} $u$ is the scalar $\det u$ such that
\[
\det{}_{\mathcal{B}}\bigl(u(x_1), \dots, u(x_n)\bigr)
= \det u \cdot \det{}_{\mathcal{B}}(x_1, \dots, x_n)
\quad \text{for all } x_i
\]
(the left side is alternating $n$-linear, hence a multiple of
$\det_\mathcal{B}$ by \cref{thm:b2:linalg:detspace}; the factor does
not depend on $\mathcal{B}$).
\end{definition}

\begin{theorem}[The determinant calculus, proved]\label{thm:b2:linalg:detrules}
\begin{enumerate}
  \item $\det(uv) = \det u\,\det v$; $\;\det(AB) = \det A \det B$.
  \item $u$ is invertible $\iff \det u \neq 0$; a family is a basis
        $\iff$ its determinant in some basis is nonzero.
  \item $\det(A^{\mathsf T}) = \det A$.
  \item Cofactor expansion along any row or column, as stated in the
        Year 1 volume, holds; similar matrices share their
        determinant.
\end{enumerate}
\end{theorem}

\begin{proof}
(1) Apply the defining relation twice:
$\det_{\mathcal B}(uv(x_i)) = \det u \cdot \det_{\mathcal
B}(v(x_i)) = \det u \det v \cdot \det_{\mathcal B}(x_i)$.

(2) If $u$ is invertible, $\det u \det u^{-1} = \det \mathrm{id} =
1 \neq 0$. If not, the images $u(e_i)$ are linked; expressing one
through the others and expanding, $\det_{\mathcal B}(u(e_i)) = 0$
(alternating kills repeated directions), so $\det u = 0$. The basis
criterion is the same statement for families.

(3) In the permutation formula, reindex each product by $j =
\sigma(i)$, i.e.\ $i = \tau(j)$ with $\tau = \sigma^{-1}$: the
factors are the same numbers in a different order, so
\[
\prod_{i=1}^{n} a_{\sigma(i),\,i} = \prod_{j=1}^{n}
a_{j,\,\tau(j)} ,
\]
and $\varepsilon(\tau) = \varepsilon(\sigma)^{-1} =
\varepsilon(\sigma)$ (values are $\pm1$; $\varepsilon$ is a
morphism). Summing over $\sigma$ is the same as summing over
$\tau$ (inversion is a bijection of $\mathfrak{S}_n$):
\[
\det A = \sum_{\tau}\varepsilon(\tau)\prod_j a_{j,\tau(j)}
= \det(A^{\mathsf T}),
\]
the last sum being the permutation formula applied to the
transposed entries $(A^{\mathsf T})_{ij} = a_{ji}$.

(4) Fix column $j$ and split $x_j = \sum_i a_{ij} e_i$ by
linearity: $\det A = \sum_i a_{ij}\, \det(\dots, e_i, \dots)$, and
moving $e_i$ to the last position ($n - i$ transpositions of rows,
$n - j$ of columns, via (3)) identifies $\det(\dots, e_i, \dots) =
(-1)^{i+j}\Delta_{ij}$ with the minor: exactly Year 1's cofactor
rule. Similarity: $\det(P^{-1}AP) = \det P^{-1}\det A \det P =
\det A$ by (1).
\end{proof}

\begin{example}[Cofactor expansion, executed]\label{ex:b2:linalg:cofactorwork}
Compute
\[
\det\begin{pmatrix}
2 & 1 & 3\\
0 & 4 & 1\\
1 & 2 & 0
\end{pmatrix}
\]
along the first column (two zeros' worth of laziness: one).
Signs follow the checkerboard $(-1)^{i+j}$:
\[
2\,\det\begin{pmatrix}4 & 1\\ 2 & 0\end{pmatrix}
- 0
+ 1\cdot\det\begin{pmatrix}1 & 3\\ 4 & 1\end{pmatrix}
= 2(0 - 2) + (1 - 12) = -15 .
\]
Cross-check by Sarrus (\cref{ex:b2:linalg:sarrus}): $0 + 1 + 0
- 12 - 0 - 4 = -15$. Strategy, not doctrine: expand along the
line with the most zeros, and when none has any, make some
first by row operations --- one round of elimination costs
less than two cofactor layers.
\end{example}

\begin{example}[A determinant by the rules]\label{ex:b2:linalg:onesmatrix}
Let $J \in \mathcal{M}_n(K)$ be the all-ones matrix and $a \in K$;
we compute $\det(aI_n + J)$ with the tools just proved. Every
column of $aI_n + J$ sums the same way: add all rows to the first
(the determinant is unchanged --- adding a multiple of one row to
another adds a repeated-direction term, killed by alternation).
The first row becomes $(a + n, a + n, \dots, a + n)$; factor out
$a + n$ by linearity in that row, then subtract the first column
from every other column: what remains is triangular with diagonal
$(1, a, \dots, a)$. Hence
\[
\det(aI_n + J) = (a + n)\,a^{\,n-1}.
\]
The closing insight: the roots $a = 0$ (multiplicity $n - 1$) and
$a = -n$ say that $J$ has eigenvalue $0$ with multiplicity $n - 1$
and eigenvalue $n$ once --- the spectrum of the rank-one matrix
$J$, one chapter early (\cref{ch:b2:reduction} will make this
systematic).
\end{example}

\begin{example}[A determinant by the permutation formula]\label{ex:b2:linalg:permexample}
For a matrix with many zeros the formula is practical by itself:
in
\[
A = \begin{pmatrix}
0 & a & 0 & 0\\
0 & 0 & b & 0\\
0 & 0 & 0 & c\\
d & 0 & 0 & 0
\end{pmatrix},
\]
the only permutation picking nonzero entries is the $4$-cycle
$\sigma = (1\,2\,3\,4)$ mapping column $1 \to$ row $4$, etc.;
$\varepsilon(\sigma) = (-1)^3 = -1$, so $\det A = -abcd$. (Check
via three column swaps to reach a diagonal matrix.)
\end{example}

\begin{example}[A Vandermonde by the product formula]\label{ex:b2:linalg:vandermondework}
For the nodes $0, 1, 2$ (used by quadrature rules like
\cref{exo:b2:linalg:4}'s), the Vandermonde determinant of
\cref{exo:b2:linalg:11} evaluates in one glance:
\[
\det\begin{pmatrix}
1 & 1 & 1\\
0 & 1 & 2\\
0 & 1 & 4
\end{pmatrix}
= (1 - 0)(2 - 0)(2 - 1) = 2 ,
\]
and by direct expansion along the first column: $1\cdot(4 - 2)
= 2$: agreement. Nonvanishing for distinct nodes is the whole
theory of interpolation in one determinant: the evaluation
forms $P \mapsto P(a_i)$ are a basis of the dual exactly when
this determinant is nonzero, i.e.\ always for distinct $a_i$
--- \cref{ex:b2:linalg:dualexamples} quantified.
\end{example}

\section{Trace, revisited}

\begin{proposition}\label{prop:b2:linalg:trace}
The trace $\operatorname{tr} \colon \mathcal{M}_n(K) \to K$ is the
unique linear form with $\operatorname{tr}(AB) =
\operatorname{tr}(BA)$ and $\operatorname{tr}(I_n) = n$ (for
$\operatorname{char} K = 0$); the trace of an endomorphism is
well defined via any matrix representation, and
\[
\operatorname{tr}(u) = \sum_{i} e_i^*\bigl(u(e_i)\bigr)
\]
in any basis --- duality writes the trace basis-freely.
\end{proposition}

\begin{proof}
$\operatorname{tr}(AB) = \operatorname{tr}(BA)$ and basis-invariance
were proved in Year 1. Uniqueness: a linear form $t$ with $t(AB) =
t(BA)$ kills every commutator $AB - BA$. We claim the commutators
span the trace-zero hyperplane, of dimension $n^2 - 1$. Two
families of commutators suffice. The multiplication rule of the
elementary matrices is $E_{ab}E_{cd} = \delta_{bc}E_{ad}$. For
$i \neq j$ it gives
\[
E_{ii}E_{ij} - E_{ij}E_{ii} = E_{ij} - 0 = E_{ij}
\]
(the second product is $E_{ij}E_{ii} = \delta_{ji}E_{ii} = 0$
since $j \neq i$): every off-diagonal $E_{ij}$ is a commutator.
And
\[
E_{ij}E_{ji} - E_{ji}E_{ij} = E_{ii} - E_{jj} .
\]
The $E_{ij}$ ($i \neq j$, $n^2 - n$ of them) together with the
$E_{11} - E_{jj}$ ($j \geq 2$, $n - 1$ of them) are $n^2 - 1$
linearly independent trace-zero matrices: they span the
hyperplane $\ker\operatorname{tr}$. So $t$ vanishes where
$\operatorname{tr}$ does and factors through it: $t =
c\operatorname{tr}$; then $t(I) = n$ forces $c = 1$. The display:
the $i$-th diagonal entry of the matrix of $u$ is precisely
$e_i^*(u(e_i))$.
\end{proof}

\begin{remark}[Common pitfalls]\label{rem:b2:linalg:pitfalls}
(i) The determinant is $n$-linear in the \emph{columns}, not
linear in the matrix: $\det(A + B) \neq \det A + \det B$ in
general, and $\det(\lambda A) = \lambda^n\det A$, not
$\lambda\det A$. (ii) Transposition reverses products:
$(vu)^{\mathsf T} = u^{\mathsf T}v^{\mathsf T}$; forgetting the
reversal wrecks every computation involving inverses. (iii) The
annihilator $F^\circ$ lives in $E^*$, not in $E$: it becomes the
familiar ``orthogonal complement'' only after an inner product
identifies $E$ with $E^*$ (\cref{ch:b2:quadratic}); no such
identification is canonical. (iv) ``Row rank equals column
rank'' does not mean row \emph{space} equals column space ---
the two live in different spaces ($K^n$ and $K^m$) and are
related through \cref{prop:b2:linalg:transposerank}, not equal.
(v) The permutation formula is a proof device: for numbers, use
row operations and cofactors (\cref{ex:b2:linalg:sarrus}).
\end{remark}

\begin{example}[The trace pairing splits the matrix space]\label{ex:b2:linalg:tracepairing}
On $\mathcal{M}_2(\R)$ with the pairing $\langle A, B\rangle =
\operatorname{tr}(AB)$ of \cref{exo:b2:linalg:9}: decompose $M =
\begin{psmallmatrix}1 & 4\\ 2 & 3\end{psmallmatrix}$ into
symmetric and antisymmetric parts,
\[
M = S + A, \qquad
S = \tfrac12(M + M^{\mathsf T}) =
\begin{pmatrix}1 & 3\\ 3 & 3\end{pmatrix},
\qquad
A = \tfrac12(M - M^{\mathsf T}) =
\begin{pmatrix}0 & 1\\ -1 & 0\end{pmatrix}.
\]
Then $\operatorname{tr}(SA) = \operatorname{tr}
\begin{psmallmatrix}-3 & 1\\ -3 & 3\end{psmallmatrix} = 0$: the
two parts are ``orthogonal'' for the trace pairing --- an
instance of the general fact (proved in the weekend problem of
this chapter) that antisymmetric matrices form exactly the
annihilator of the symmetric ones. Duality sees the
decomposition $\mathcal{M}_n = \mathcal{S}_n \oplus
\mathcal{A}_n$ before any inner product is chosen.
\end{example}

\begin{remark}[Perspectives within this volume]\label{rem:b2:linalg:perspectives}
Watch the three constructions of this chapter change costume
ahead. The \emph{transpose} returns in
\cref{ch:b2:reduction}: $u$ and $u^{\mathsf T}$ share
eigenvalues with equal geometric multiplicities (this chapter's
weekend problem, question 15), which is why row and column
analyses of a matrix never disagree. The \emph{determinant}
becomes a function of a parameter in \cref{ch:b2:reduction}
($\chi_u(X) = \det(X\,\mathrm{id} - u)$) and a Jacobian in
\cref{ch:b2:multint}, where its multilinearity turns into the
change-of-variables factor. The \emph{trace} seeds the
similarity invariants: it is the second coefficient of
$\chi_u$, the sum of eigenvalues, and eventually the integral
of the diagonal in \cref{ch:b2:fourier}-style identities. One
linear-algebra chapter, three long shadows.
\end{remark}

\begin{remark}[Where this chapter is used]
The dual space is not an abstraction for its own sake: annihilators
and transposes run the solvability theory of linear systems (this
chapter's weekend problem proves the finite-dimensional Fredholm
alternative from them), nondegenerate pairings reappear as the
polar form in \cref{ch:b2:quadratic} and the adjoint in
\cref{ch:b2:hermitian}, and the determinant built here powers the
whole of \cref{ch:b2:reduction}. In the Year 3 volume the same
duality, transported to infinite dimension, becomes the
Riesz representation theorem and Fredholm theory on Hilbert
spaces --- with compactness replacing the dimension counts used
here.
\end{remark}

\section{Exercises}

\begin{exercise}[$\star$]\label{exo:b2:linalg:1}
In $\R^3$, let $\varphi_1(x,y,z) = x + y$, $\varphi_2 = y + z$,
$\varphi_3 = x + z$. Prove that $(\varphi_1, \varphi_2, \varphi_3)$
is a basis of $(\R^3)^*$ and find the basis of $\R^3$ of which it is
the dual.
\end{exercise}

\begin{exercise}[$\star$]\label{exo:b2:linalg:2}
Compute by the permutation formula the determinants of
\[
\begin{pmatrix} 0 & 0 & a\\ 0 & b & 0\\ c & 0 & 0 \end{pmatrix},
\qquad
\begin{pmatrix}
a & b & 0 & 0\\
c & d & 0 & 0\\
0 & 0 & e & f\\
0 & 0 & g & h
\end{pmatrix},
\]
and state the block-diagonal rule the second suggests.
\end{exercise}

\begin{exercise}[$\star$]\label{exo:b2:linalg:3}
Let $F = \{(x,y,z,t) \in \R^4 : x + y = z + t \text{ and } x = 2y\}$.
Give a basis of $F^\circ$ and check
\cref{thm:b2:linalg:annihilator} on dimensions.
\end{exercise}

\begin{exercise}[$\star\star$]\label{exo:b2:linalg:4}
Let $a_0, \dots, a_n$ be distinct points of $K$ and $\varphi_i \colon
P \mapsto P(a_i)$ on $K_n[X]$. Prove that $(\varphi_0, \dots,
\varphi_n)$ is a basis of $K_n[X]^*$, identify its pre-dual basis,
and expand the form $P \mapsto \int_0^1 P(t)\,\dd t$ (for $K = \R$,
$n = 2$, $a_i = 0, \frac12, 1$) in this basis --- recognizing
Simpson's rule.
\end{exercise}

\begin{exercise}[$\star\star$]\label{exo:b2:linalg:5}
Let $u \in \mathcal{L}(E)$ with $\dim E = n$ and $\operatorname{rk} u
= 1$. Prove that $u = \varphi(\cdot)\, a$ for a vector $a$ and a
form $\varphi$; that $\operatorname{tr} u = \varphi(a)$; and that
$u^2 = (\operatorname{tr} u)\, u$. Deduce $\det(I + u) = 1 +
\operatorname{tr} u$.
\end{exercise}

\begin{exercise}[$\star\star$]\label{exo:b2:linalg:6}
Prove that every hyperplane of $\mathcal{M}_n(K)$ ($n \geq 2$)
contains an invertible matrix. \emph{Hint: a hyperplane is $\{M :
\operatorname{tr}(AM) = 0\}$ for some $A \neq 0$
(\cref{exo:b2:linalg:9}). If $A$ is scalar, exhibit an invertible
matrix of zero trace; otherwise, find an invertible $M$ making
$AM$ have zero diagonal --- a permutation-like matrix does it.}
\end{exercise}

\begin{exercise}[$\star\star$]\label{exo:b2:linalg:7}
(Derivative of the determinant) For $A \in \mathcal{M}_n(\R)$, prove
from multilinearity that
\[
\frac{\dd}{\dd t}\Big|_{t=0} \det(I_n + tA) = \operatorname{tr} A ,
\]
and deduce $\det(\eu^{tA}) = \eu^{t\operatorname{tr} A}$ assuming
the differentiability of $t \mapsto \det(\eu^{tA})$ and the group
property $\eu^{(s+t)A} = \eu^{sA}\eu^{tA}$ (established in
\cref{ch:b2:diffeq}).
\end{exercise}

\begin{exercise}[$\star\star$]\label{exo:b2:linalg:8}
(Circulant, $3 \times 3$) Let $j = \eu^{2\iu\pi/3}$ and
\[
C = \begin{pmatrix}
a & b & c\\
c & a & b\\
b & c & a
\end{pmatrix} \in \mathcal{M}_3(\C).
\]
Verify that the columns of the Vandermonde matrix of $1, j, j^2$ are
eigenvectors of $C$, and deduce
\[
\det C = (a + b + c)(a + bj + cj^2)(a + bj^2 + cj).
\]
\end{exercise}

\begin{exercise}[$\star\star\star$]\label{exo:b2:linalg:9}
Prove that every linear form $t$ on $\mathcal{M}_n(K)$ is $M \mapsto
\operatorname{tr}(AM)$ for a unique $A$: the map $A \mapsto
\operatorname{tr}(A\,\cdot)$ is an isomorphism from
$\mathcal{M}_n(K)$ onto its dual. Deduce the uniqueness statement of
\cref{prop:b2:linalg:trace} again.
\end{exercise}

\begin{exercise}[$\star\star\star$]\label{exo:b2:linalg:10}
Let $u, v \in \mathcal{L}(E)$ with $u \circ v - v \circ u = u$.
Prove that $u$ is nilpotent. \emph{Hint: show
$\operatorname{tr}(u^k) = 0$ for all $k \geq 1$ (compute $u^k v -
v u^k$ by induction), then use the following fact, to be proved via
Newton's identities or by induction on the dimension: an
endomorphism of a $\C$-vector space all of whose powers have zero
trace is nilpotent. Work over $\C$.}
\end{exercise}

\begin{exercise}[$\star\star$]\label{exo:b2:linalg:11}
(Vandermonde) For $a_0, \dots, a_n \in K$, prove
\[
\det\begin{pmatrix}
1 & 1 & \cdots & 1\\
a_0 & a_1 & \cdots & a_n\\
\vdots & \vdots & & \vdots\\
a_0^n & a_1^n & \cdots & a_n^n
\end{pmatrix}
= \prod_{0 \leq i < j \leq n} (a_j - a_i).
\]
\emph{(View the determinant as a polynomial in $a_n$: identify its
degree, its roots, and its leading coefficient; induct.)}
\end{exercise}

\begin{exercise}[$\star\star\star$]\label{exo:b2:linalg:12}
Let $A, B, C, D \in \mathcal{M}_n(K)$ with $K$ infinite, and
suppose $CD = DC$. Prove that
\[
\det\begin{pmatrix} A & B\\ C & D\end{pmatrix}
= \det(AD - BC).
\]
\emph{(Treat first $D$ invertible, multiplying on the right by
$\begin{psmallmatrix} I & 0\\ -D^{-1}C & I\end{psmallmatrix}$;
then replace $D$ by $D + tI$ and compare two polynomials in $t$.)}
\end{exercise}

\section{Problem: The Fredholm Alternative}

When does the linear system $u(x) = b$ have a solution? The
complete answer is a duality statement: \emph{exactly when $b$ is
annihilated by every linear form that annihilates the image of
$u$} --- and those forms are computable, being the kernel of the
transpose. This weekend problem builds the full dictionary of
finite-dimensional duality (factorization of forms, biduality,
annihilator calculus, the transpose), proves the
finite-dimensional \emph{Fredholm alternative}, and closes with
the trace form and a characterization: the trace is the only
linear invariant of similarity. Throughout, $E$ and $F$ are
finite-dimensional $K$-vector spaces, $n = \dim E$.

\begin{problem}[{Weekend problem --- duality in finite dimension
and the Fredholm alternative}]
\label{pb:b2:linalg:1}
Notation: for $S \subseteq E^*$, the \emph{pre-annihilator} is
$S_\circ = \{x \in E : \varphi(x) = 0 \text{ for all } \varphi
\in S\}$; annihilators $F^\circ$ and transposes $u^{\mathsf T}$
are those of \cref{def:b2:linalg:annihilator} and
\cref{def:b2:linalg:transpose}.

\medskip
\noindent\textbf{Part I --- The factorization lemma.} Let
$\varphi_1, \dots, \varphi_p, \varphi \in E^*$.
\begin{enumerate}
  \item Let $\Phi \colon E \to K^p$, $x \mapsto (\varphi_1(x),
        \dots, \varphi_p(x))$. Identify $\ker\Phi$, show
        $\Phi^{\mathsf T}$ maps the coordinate forms of $K^p$ to
        the $\varphi_i$, and deduce
        \[
        \dim \bigl(\ker\varphi_1 \cap \dots \cap
        \ker\varphi_p\bigr) = n - \dim
        \operatorname{Vect}(\varphi_1, \dots, \varphi_p).
        \]
  \item (Factorization lemma) Prove the equivalence:
        \[
        \varphi \in \operatorname{Vect}(\varphi_1, \dots,
        \varphi_p)
        \iff
        \ker\varphi_1 \cap \dots \cap \ker\varphi_p \subseteq
        \ker\varphi .
        \]
  \item Deduce: $(\varphi_1, \dots, \varphi_p)$ is free iff
        $\bigcap_i \ker\varphi_i$ has dimension $n - p$; and a
        subspace of codimension $p$ is an intersection of $p$
        hyperplanes, never fewer.
  \item In $\R^4$, let $\varphi_1 = x + y - z$, $\varphi_2 = y +
        z - t$, $\psi = x + 2y - t$ and $\psi' = x + y + t$.
        Decide, by the factorization lemma, whether $\psi$ and
        $\psi'$ belong to $\operatorname{Vect}(\varphi_1,
        \varphi_2)$.
  \item On $E = \R_2[X]$, show that $\psi_0 \colon P \mapsto
        P(0)$, $\psi_1 \colon P \mapsto P(1)$, $\psi_2 \colon P
        \mapsto \int_0^1 P(t)\dd t$ form a basis of $E^*$,
        compute the basis $(P_0, P_1, P_2)$ of $E$ of which it is
        the dual, and find the unique $P \in \R_2[X]$ with $P(0)
        = 1$, $P(1) = 2$, $\int_0^1 P = \frac32$.
\end{enumerate}

\medskip
\noindent\textbf{Part II --- Biduality and the annihilator
calculus.}
\begin{enumerate}[resume]
  \item Show that the \emph{evaluation map} $J \colon E \to
        E^{**}$, $J(x)(\varphi) = \varphi(x)$, is linear and
        injective, hence an isomorphism in finite dimension.
  \item (Double annihilator) Show $J(F) = F^{\circ\circ} :=
        (F^\circ)^\circ$ for every subspace $F \subseteq E$: under
        the identification $J$, the annihilator of the annihilator
        is the subspace itself.
  \item Prove the annihilator calculus: $(F + G)^\circ = F^\circ
        \cap G^\circ$ and $(F \cap G)^\circ = F^\circ + G^\circ$.
  \item Deduce (and reprove directly): two nonzero forms with the
        same kernel are proportional.
  \item (Antedual basis) Show that for every basis $(\varphi_1,
        \dots, \varphi_n)$ of $E^*$ there is a unique basis
        $(u_1, \dots, u_n)$ of $E$ with $\varphi_i(u_j) =
        \delta_{ij}$.
\end{enumerate}

\medskip
\noindent\textbf{Part III --- The transpose calculus.}
\begin{enumerate}[resume]
  \item Show that $u \mapsto u^{\mathsf T}$ is a linear bijection
        from $\mathcal{L}(E, F)$ onto $\mathcal{L}(F^*, E^*)$,
        and that $(u^{-1})^{\mathsf T} = (u^{\mathsf T})^{-1}$
        when $u$ is invertible.
  \item (Naturality) Show that $u^{\mathsf T\mathsf T} \circ J_E
        = J_F \circ u$: under the evaluation isomorphisms, the
        double transpose \emph{is} $u$.
  \item Show: $u$ is surjective iff $u^{\mathsf T}$ is injective;
        $u$ is injective iff $u^{\mathsf T}$ is surjective.
  \item For $u \in \mathcal{L}(E)$: a subspace $F$ is stable
        under $u$ if and only if $F^\circ$ is stable under
        $u^{\mathsf T}$.
  \item Show that $\ker(u^{\mathsf T} - \lambda\,
        \mathrm{id}_{E^*}) = \bigl(\operatorname{im}(u - \lambda\,
        \mathrm{id}_E)\bigr)^\circ$, and deduce that $u$ and
        $u^{\mathsf T}$ have the same eigenvalues with the same
        geometric multiplicities.
\end{enumerate}

\medskip
\noindent\textbf{Part IV --- The Fredholm alternative.}
\begin{enumerate}[resume]
  \item Prove that $\operatorname{im} u = (\ker u^{\mathsf
        T})_\circ$ for $u \in \mathcal{L}(E, F)$, and deduce the
        \emph{Fredholm alternative} in finite dimension: the
        equation $u(x) = b$ has a solution if and only if every
        $\psi \in F^*$ with $u^{\mathsf T}\psi = 0$ satisfies
        $\psi(b) = 0$.
  \item Matrix form: for $A \in \mathcal{M}_{m,n}(K)$ and $b \in
        K^m$, exactly one of the following holds: (i) $Ax = b$
        has a solution; (ii) there is $y \in K^m$ with
        $A^{\mathsf T}y = 0$ and $y^{\mathsf T}b = 1$. Prove both
        the ``at most one'' and the ``at least one''.
  \item Find all $b \in \R^3$ for which the system
        \[
        x + y = b_1, \qquad y + z = b_2, \qquad x + 2y + z = b_3
        \]
        has a solution, by computing the kernel of the transposed
        matrix.
  \item (A discrete Neumann problem) On $E = \R^n$ ($n \geq 3$),
        define $L$ by $(Lx)_k = x_k - \frac12(x_{k-1} +
        x_{k+1})$, indices modulo $n$. Show $L^{\mathsf T} = L$
        (canonical identifications), show $\ker L$ is the line of
        constant vectors \emph{(look at a maximal coordinate)},
        and conclude: $Lx = b$ is solvable iff $\sum_k b_k = 0$.
\end{enumerate}

\medskip
\noindent\textbf{Part V --- The trace form and the invariance
theorem.} Recall from \cref{exo:b2:linalg:9} that $A \mapsto
\operatorname{tr}(A\,\cdot)$ identifies $\mathcal{M}_n(K)$ with
its dual. Assume $\operatorname{char} K = 0$ (e.g.\ $K = \Q, \R,
\C$).
\begin{enumerate}[resume]
  \item Under this identification, show that the annihilator of
        the subspace $\mathcal{S}_n$ of symmetric matrices is the
        subspace $\mathcal{A}_n$ of antisymmetric matrices, and
        conversely.
  \item Show that the annihilator of the hyperplane
        $\mathfrak{sl}_n = \{M : \operatorname{tr} M = 0\}$ is
        the line $K I_n$; equivalently, a linear form vanishing
        on all trace-zero matrices is a multiple of the trace.
  \item Show that every matrix of $\mathcal{M}_n(K)$ is the sum
        of two invertible matrices.
  \item (The trace is the only linear similarity invariant) Let
        $t$ be a linear form on $\mathcal{M}_n(K)$ with
        $t(PMP^{-1}) = t(M)$ for every $M$ and every invertible
        $P$. Show first $t(PX) = t(XP)$ for $P$ invertible, then
        $t(BX) = t(XB)$ for \emph{all} $B$, and conclude $t = c
        \operatorname{tr}$ for some $c \in K$.
  \item Show that $\operatorname{rk} u \leq r$ if and only if $u$
        is a sum of $r$ maps of rank $\leq 1$, i.e.\ $u =
        \sum_{i=1}^{r} \psi_i(\cdot)\,f_i$ with $\psi_i \in E^*$,
        $f_i \in F$; deduce $\operatorname{rk}(u + v) \leq
        \operatorname{rk} u + \operatorname{rk} v$.
  \item (Synthesis) Draw up the dictionary proved in this
        problem: subspaces versus annihilators, sums versus
        intersections, maps versus transposes, solvability versus
        orthogonality to the transposed kernel, trace versus
        similarity. For each entry, cite the question that proved
        it, and state in one sentence what replaces the dimension
        counts when dimension becomes infinite (the Year 3 volume
        makes this precise on Hilbert spaces).
\end{enumerate}
\end{problem}
